\chapter{BDI Agent}
\label{ch:bdi}
Both agents are built on the same Belief--Desire--Intention reasoning core: each
constructs a \texttt{BeliefBase} and runs an \texttt{IntentionEngine}. The physical
executor uses it directly to move, pick up and deliver; the LLM coordinator
(\cref{ch:llm}) reuses the same structure and extends it with
natural-language reasoning.

The agent maintains a model of the world (\emph{beliefs}), evaluates which goals
are worth pursuing (\emph{desires}), commits to one (\emph{intention}), and
executes it incrementally. A control loop, driven by the game loop, runs each cycle
of the schema below at a tick.

\begin{center}
\begin{tikzpicture}[
  node distance=2mm and 4mm,
  box/.style={draw, rounded corners=2pt, font=\small, inner sep=3pt},
  >={Stealth[round]}
]
  \node[box] (perceive)           {perceive};
  \node[box, right=of perceive]   (revise)     {revise beliefs};
  \node[box, right=of revise]     (deliberate) {deliberate};
  \node[box, right=of deliberate] (commit)     {commit to plan};
  \node[box, right=of commit]     (act)        {act};
  \node[font=\scriptsize, below=3mm] at ($(deliberate.east)!0.5!(commit.west)$) {(throttled)};

  \draw[->] (perceive)   -- (revise);
  \draw[->] (revise)     -- (deliberate);
  \draw[->] (deliberate) -- (commit);
  \draw[->] (commit)     -- (act);

  \draw[->] (act.south) -- ++(0,-5mm)
    -| node[below, pos=0.25, font=\scriptsize]{world} (revise.south);
\end{tikzpicture}
\end{center}

%\[
%\text{Beliefs} \;\rightarrow\; \text{Goal} \;\rightarrow\; \text{Plan}
%\;\rightarrow\; \text{Action} \;\rightarrow\; \text{Feedback}
%\;\rightarrow\; \text{Beliefs}
%\]

\section{Sensing and Belief Revision}
\label{sec:sensing-beliefs}
Sensing happens at each movement and whenever something new enters the observation 
zone. A raw perception frame from the server reports the agent's own state, the
visible map configuration, and the nearby agents, crates and parcels, but only
within the observation radius, so it is partial and momentary. The agent never
reasons on this raw stream directly: every frame is merged into a persistent
\texttt{BeliefBase} that combines fresh observations with previously acquired
knowledge, allowing the agent to reason about entities currently out of sight.\\

Revision performs three non-trivial inferences:
\begin{itemize}
  \item \textbf{Spatial memory.} Parcels and crates that leave the field of view
        are retained rather than forgotten.
  \item \textbf{Negative inference.} If a remembered object should be visible, since it falls inside
        the current radius, but it is absent from the frame, the agent
        concludes it was collected, decayed, or moved, and prunes it.
  \item \textbf{Decay simulation.} Carried and remembered parcel rewards are aged
        locally over elapsed time, so value estimates stay accurate between
        sightings.
\end{itemize}

The belief base also holds state never delivered by perception: the agent's own
identity (\texttt{me}), the carried inventory, peers, locked targets, active
cooperative contracts, and the \texttt{policyRules} issued by the LLM coordinator.
A game loop then calls the \texttt{IntentionEngine} to reason and act on the current
beliefs, while an \texttt{onMsg} handler processes Admin prompts (for the LLM agent)
and the P2P coordination messages exchanged between the two agents.

\section{Navigation Graph (Adjacency)}
\label{sec:adjacency}
There is no adjacency map sent by the server: the initial state (\texttt{onMap})
only provides a flat list of tiles, each with $x$, $y$ and a type. Adjacency, which is the
edges of the navigation graph describing which tile-to-tile moves are legal, is
derived on demand from these tile codes by \texttt{MapRepresentation.js}.

Two methods implement this. \texttt{isAdjacent(from, to)} checks whether a single
step is allowed: the target tile must be exactly one step away, must be walkable,
and the move must not violate the one-way arrow gate tiles (a tile cannot be entered
against its arrow direction). \texttt{getNeighbors(pos)} returns all legal neighbors
by applying \texttt{isAdjacent} to the four cardinal directions.

Because of the one-way gates, adjacency is directed and not symmetric: a move from
$A$ to $B$ may be legal while the reverse is not. For this reason no adjacency table
is precomputed; the graph is implicit and evaluated lazily. Pathfinding
(\texttt{src/mapping/Pathfinding.js}) expands it by repeatedly calling
\texttt{getNeighbors()} during the search, automatically respecting walls and gates.

\section{Desires and Goal Selection}
\label{sec:desires-goals}
Deliberation is throttled to stay stable and avoid redundant pathfinding: goal
re-evaluation runs every \texttt{GOAL\_EVAL\_INTERVAL} (12 ticks, about 200ms), but
fires immediately when there is no active plan or when the carried inventory
changes. When it runs, \texttt{selectBestGoal(beliefs, ...)}
(\texttt{GoalSelector.js}) reads the beliefs and returns a single goal descriptor
\texttt{\{type, targetId, x, y\}}, chosen by a strict priority cascade:

\begin{enumerate}
  \item admin direct commands (\texttt{admin\_move},
        \texttt{admin\_pickup}, \texttt{admin\_deliver});
  \item active cooperative contracts with the partner
        (\texttt{rendezvous}, \texttt{handoff}, \texttt{relay});
  \item an economic, utility-based evaluation over the perceived parcels.
\end{enumerate}

The economic evaluation is the core decision and is \emph{expiration-aware}. For
each candidate parcel it estimates the value of the entire trip using A\textsuperscript{*}
travel-time estimates, discounting the raw reward by the decay accrued over that
time and adjusting it by the coordinator's policy multipliers and bonuses:
%
\begin{equation}
  \mathrm{value} = \bigl(\mathrm{reward} - r_{\text{decay}}\cdot t_{\text{trip}}\bigr)
                   \cdot m_{\text{policy}} + b_{\text{policy}}, \qquad
  \mathrm{utility} = \frac{\mathrm{value}}{t_{\text{trip}}}.
\end{equation}
%
Because travel time is real, a closer modest parcel can outscore a distant high-reward
one. When already carrying cargo, the agent also weighs delivering now against detouring
for one more pickup. This trade-off is bounded by three factors: the carrying capacity,
the decay risk to the cargo already held, and any required-stack-size incentive, which
can make it worth delaying delivery to assemble a larger batch. When no option scores
positively, deliberation degrades gracefully: it delivers what is held, anchors near a
relay tile, patrols the parcel spawn zones, and finally wanders.

\subsection{Delivery Optimization}
\label{sub:delivery-opt}
Both the utility estimates above and the delivery step itself defer to a dedicated
optimizer \\
(\texttt{DeliveryOptimizer.js}): on reaching a delivery tile the agent rarely
just drops everything, since policy rules may bound which rewards score or require a
minimum stack size. \texttt{optimizeDeliveryStack} jointly chooses \emph{which subset}
of cargo to deliver and \emph{how long to wait} first (candidate wait times are taken
from the policy reward boundaries, the instants a parcel crosses an allowed band).
Each (subset, wait) pair is scored as
%
\begin{equation}
  \mathrm{score} = R_{\text{delivered}} + \tfrac{1}{2}\,V_{\text{kept}}
                   - \lambda\, t_{\text{wait}},
\end{equation}
%
balancing the reward actually delivered against the discounted value of cargo kept back
and a waiting penalty; parcels that can never yield positive value are discarded. The
same routine is reused during goal selection, keeping valuation and execution consistent.

\section{Intention Commitment and Plans}
\label{sec:intention-plans}
The agent is deliberately reluctant to abandon an intention: a newly selected goal
replaces the running one only when \texttt{shouldPreemptActivePlan(...)}
(\texttt{PreemptionManager.js}) judges the switch worthwhile, which prevents
oscillation between comparable options. The exception is a path blocked by a crate,
which suspends the active plan and injects a PDDL \texttt{clear\_corridor} plan
instead (\cref{ch:pddl}).

On commit, \texttt{instantiatePlanRecipe(goal)} maps the goal type to a
generator-based plan recipe (e.g. \texttt{pickup} $\to$ \texttt{CollectAndDeliver},
\texttt{deliver} $\to$ \texttt{\_deliverRecipe}, plus the admin and cooperation
recipes). These generators yield one primitive step at a time and can be suspended
and resumed via a suspension stack, so a corridor-clearing plan can interrupt a
delivery and the delivery resumes afterwards. Suspended plans that have gone stale, since
their target disappeared or became obsolete, are discarded rather than resumed.

\section{Execution}
\label{sec:exec}
Each cycle advances the active plan by one step: (\texttt{move},
\texttt{pickup}, \texttt{putdown}, \texttt{wait}, \texttt{say}), handed to
\texttt{dispatchAction} (\texttt{ActionDispatcher.js}) which emits it on the socket.
The plan is told whether the previous step succeeded, so it can react to failures
rather than continue blindly.

Two behaviours sit outside the normal deliberation cycle. The first is an
\emph{opportunistic pickup}: if the agent finds itself standing on a free parcel and
is still below its carrying capacity (starting at infinity), it grabs the parcel immediately, before any
goal evaluation (deliberation). This is done to avoid leaving reachable parcels on the floor and wasting potential value.

The second is \emph{collision avoidance} between the two teammates, designed so they
never deadlock over a tile during testing. Before moving, each agent uses the partner's broadcast
position and planned next step to detect a possible conflict where both aim for the same tile (see
\cref{sec:agent-comm}).
When this happens one agent must yield, and this is decided deterministically by comparing agent identifiers,
so both reach the same conclusion independently without any negotiation(higher ID yields). The yeilding agent 
pauses \SI{150}{\milli\second} to let the other pass. If a move is still rejected,
the agent waits one tick (\SI{100}{\milli\second}) and retries, and after two repeated
collisions on the same step it marks that tile as blocked and re-plans a route around it.
The block is cleared after a short timeout so the tile stays reusable.
When a plan finishes, the engine resumes any suspended plan or clears the goal and
re-deliberates on the next tick.



%If a collision occurs anyway, the agent waits one tick and retries; if the same tile
%keeps rejecting it across repeated attempts, it marks the tile as blocked and re-plans
%a route around it, clearing the mark again after a short timeout so the tile stays reusable.